\begin{ex}
Tính: $9 + 7$.
\shortans{16}
\loigiai{
Ta có $9 + 7 = 16$.
}
\end{ex}

\begin{ex}
Tính: $16 - 9$.
\shortans{7}
\loigiai{
Ta có $16 - 9 = 7$.
}
\end{ex}

\begin{ex}
Tính: $8 + 6$.
\shortans{14}
\loigiai{
Ta có $8 + 6 = 14$.
}
\end{ex}

\begin{ex}
Tính: $14 - 8$.
\shortans{6}
\loigiai{
Ta có $14 - 8 = 6$.
}
\end{ex}

\begin{ex}
Tính: $7 + 5$.
\shortans{12}
\loigiai{
Ta có $7 + 5 = 12$.
}
\end{ex}

\begin{ex}
Tính: $12 - 7$.
\shortans{5}
\loigiai{
Ta có $12 - 7 = 5$.
}
\end{ex}

\begin{ex}
Tính giá trị biểu thức: $9 + 5 - 6$.
\shortans{8}
\loigiai{
Ta có $9 + 5 = 14$, rồi $14 - 6 = 8$.
}
\end{ex}

\begin{ex}
Tính giá trị biểu thức: $13 - 7 + 8$.
\shortans{14}
\loigiai{
Ta có $13 - 7 = 6$, rồi $6 + 8 = 14$.
}
\end{ex}

\begin{ex}
Tính giá trị biểu thức: $15 - 8 + 9$.
\shortans{16}
\loigiai{
Ta có $15 - 8 = 7$, rồi $7 + 9 = 16$.
}
\end{ex}

\begin{ex}
Tính giá trị biểu thức: $8 + 8 - 9$.
\shortans{7}
\loigiai{
Ta có $8 + 8 = 16$, rồi $16 - 9 = 7$.
}
\end{ex}

\begin{ex}
Điền số thích hợp: $7 + 8 = 9 + \dots$
\shortans{6}
\loigiai{
Vế trái: $7 + 8 = 15$. Số cần điền: $15 - 9 = 6$.
}
\end{ex}

\begin{ex}
Điền số thích hợp: $16 - 7 = 14 - \dots$
\shortans{5}
\loigiai{
Vế trái: $16 - 7 = 9$. Số cần điền: $14 - 9 = 5$.
}
\end{ex}

\begin{ex}
Điền số thích hợp: $13 - 6 = \dots - 8$.
\shortans{15}
\loigiai{
Vế trái: $13 - 6 = 7$. Số cần điền: $7 + 8 = 15$.
}
\end{ex}

\begin{ex}
Điền số thích hợp: $8 + 5 = \dots + 4$.
\shortans{9}
\loigiai{
Vế trái: $8 + 5 = 13$. Số cần điền: $13 - 4 = 9$.
}
\end{ex}

\begin{ex}
Một đàn gà có 15 con, trong đó có 7 con gà trống. Hỏi đàn gà có bao nhiêu con gà mái?
\shortans{8}
\loigiai{
Số gà mái: $15 - 7 = 8$ con.
}
\end{ex}

\begin{ex}
Lan hái được 8 bông hoa hồng, Mai hái được nhiều hơn Lan 5 bông hoa hồng. Hỏi Mai hái được bao nhiêu bông hoa hồng?
\shortans{13}
\loigiai{
Mai hái được: $8 + 5 = 13$ bông hoa.
}
\end{ex}

\begin{ex}
Hùng có 17 viên bi, Hùng cho Dũng 9 viên bi. Hùng còn lại bao nhiêu viên bi?
\shortans{8}
\loigiai{
Hùng còn lại: $17 - 9 = 8$ viên bi.
}
\end{ex}

\begin{ex}
Bình có 14 cái nhãn vở, Minh có ít hơn Bình 6 cái nhãn vở. Hỏi Minh có bao nhiêu cái nhãn vở?
\shortans{8}
\loigiai{
Minh có: $14 - 6 = 8$ cái nhãn vở.
}
\end{ex}

\begin{ex}
Cả hai bạn Lan và Mai có 16 cái kẹo, biết Lan có 9 cái kẹo. Hỏi Mai có bao nhiêu cái kẹo?
\shortans{7}
\loigiai{
Mai có: $16 - 9 = 7$ cái kẹo.
}
\end{ex}

\begin{ex}
Hiệu của 18 và số lớn nhất có một chữ số là bao nhiêu?
\shortans{9}
\loigiai{
Số lớn nhất có 1 chữ số là 9. Hiệu: $18 - 9 = 9$.
}
\end{ex}

\begin{ex}
Tổng của 8 và số liền sau của 6 là bao nhiêu?
\shortans{15}
\loigiai{
Số liền sau của 6 là 7. Tổng: $8 + 7 = 15$.
}
\end{ex}

\begin{ex}
Hiệu của 14 và số liền trước của 8 là bao nhiêu?
\shortans{7}
\loigiai{
Số liền trước của 8 là 7. Hiệu: $14 - 7 = 7$.
}
\end{ex}

\begin{ex}
Bác thợ rèn thanh sắt dài 15 cm, bác cắt bớt một đoạn dài 8 cm. Đoạn sắt còn lại dài bao nhiêu xăng-ti-mét?
\shortans{7}
\loigiai{
Còn lại: $15 - 8 = 7$ cm.
}
\end{ex}

\begin{ex}
Một cuộn len dài 16 dm, dùng đan khăn hết 7 dm. Cuộn len còn lại dài bao nhiêu đề-xi-mét?
\shortans{9}
\loigiai{
Còn lại: $16 - 7 = 9$ dm.
}
\end{ex}

\begin{ex}
Điền số thích hợp vào chỗ chấm: $17 - 9 = 2 + \dots$
\shortans{6}
\loigiai{
Vế trái: $17 - 9 = 8$. Số cần điền: $8 - 2 = 6$.
}
\end{ex}
